% =========================================================
% Chapter: Conclusion and Future Work
% Purpose:
% - Summarize what was achieved in both phases and why the findings matter
% - Distill key insights and acknowledge limitations
% - Lay out concrete, actionable future directions
% Navigation:
% - Sections: Contributions → Key Findings → Limitations → Future Work → Closing Remarks
% Tips:
% - Avoid introducing new results; keep this chapter concise and forward-looking
% =========================================================
\chapter{Conclusion and Future Work}
\label{ch:conclusion-future}

This thesis addressed the challenge of detecting increasingly realistic deepfakes by pursuing two complementary research directions. The first focused on a hybrid video-only detector that integrates spatial and temporal reasoning to identify visual forgeries. The second advanced multimodal detection by explicitly modeling audio--visual synchronicity via cross-modal attention and bidirectional temporal modeling. Together, these contributions demonstrate that combining strong spatial encoders with sequence-aware mechanisms and synchronization-focused cross-modal fusion yields robust performance across diverse manipulation types, datasets, and operating conditions.

% Context: Concise recap of what the thesis built and evaluated across both phases
\section{Summary of Contributions}
The main scientific and engineering contributions are summarized below:
\begin{itemize}
    \item \textbf{Hybrid Visual Detector (Phase I):} A video-only architecture combining ResNet50 for spatial encoding with Bidirectional GRUs and Multi-Head Attention for temporal focus. Trained on an aggregated corpus (FaceForensics++, Celeb-DF, DFDC), the model achieved strong generalization with high AUC and F1, surpassing baselines that emphasize solely spatial cues.
    \item \textbf{Synchronicity-Aware Multimodal Detector (Phase II):} A dual-stream audio--visual framework that processes video and audio with specialized encoders and per-stream Bi-GRUs, fused via cross-modal attention to detect viseme--phoneme misalignments. On a consolidated dataset (FakeAVCeleb and a stratified AV-Deepfake1M++ subset), the model achieved 92.5\% accuracy and 0.94 AUC-ROC, significantly outperforming unimodal and naive fusion baselines.
    \item \textbf{Robustness and Fairness Evaluation:} Systematic assessments under lossy compression and stratified demographic analyses indicated improved resilience relative to unimodal baselines and limited performance variance across groups when trained with balanced subsets.
    \item \textbf{Unified Experimental Protocols:} Consistent preprocessing, augmentation, and evaluation pipelines for visual-only and multimodal settings, facilitating reproducibility and fair comparisons across phases.
\end{itemize}

% Context: High-level takeaways from empirical studies; what worked and why
\section{Key Findings}
The empirical studies yield several insights:
\begin{itemize}
    \item \textbf{Temporal modeling matters:} Explicit sequence modeling via Bi-GRUs improves detection of temporal artifacts that are not captured by single-frame CNNs.
    \item \textbf{Cross-modal attention is critical:} Learned audio--visual alignment provides a robust cue for exposing short-duration asynchronies that persist even as spatial and spectral artifacts diminish.
    \item \textbf{Diverse training improves generalization:} Aggregated datasets and consistent augmentation reduce overfitting to dataset-specific artifacts and improve cross-dataset performance.
    \item \textbf{Operational robustness:} Multimodal synchrony remains informative under moderate compression, improving reliability in realistic content pipelines relative to unimodal detectors.
\end{itemize}

% Context: Constraints that bound the applicability of the current approach
\section{Limitations}
While the proposed frameworks mark progress, several limitations remain:
\begin{itemize}
    \item \textbf{Severe degradations:} Under extreme compression or very low resolutions, spatial and spectral cues erode, reducing the effectiveness of both phases.
    \item \textbf{Tightly synchronized forgeries:} Advanced lip-sync and high-quality TTS/VC pipelines can closely align visemes and phonemes, diminishing synchronicity cues and requiring longer-range or token-aware modeling.
    \item \textbf{Content variability:} Occlusions, off-axis poses, or non-speech mouth motions can confound cross-modal alignment and require robust face tracking and speech activity detection.
    \item \textbf{Resource demands:} Dual-stream recurrent and attention modules increase inference costs relative to lightweight single-frame models, complicating real-time deployment without compression or distillation.
\end{itemize}

% Context: Concrete next steps to extend capability, generalization, and deployability
\section{Future Work}
Several directions can extend the present work:
\begin{enumerate}
    \item \textbf{Transformer-Based Long-Range Modeling:} Replace or augment Bi-GRUs with temporal transformers to capture long-range dependencies and improve modeling of coarticulation and prosody across longer contexts.
    \item \textbf{Phoneme/Viseme-Aware Supervision:} Incorporate auxiliary alignment objectives (e.g., CTC or contrastive losses) using phoneme/viseme annotations to enhance synchronization sensitivity and interpretability.
    \item \textbf{Resolution-Invariant and Frequency-Robust Features:} Explore multi-scale encoders and frequency-aware training to increase resilience under low-resolution and high-compression conditions.
    \item \textbf{Adversarial and Evasion Robustness:} Integrate adversarial training and red-teaming to harden the detector against evasion strategies and emerging synthesis techniques.
    \item \textbf{Cross-Lingual and Cross-Cultural Generalization:} Evaluate synchronization cues across languages with diverse phonetic inventories and cultural speaking styles to ensure broad applicability.
    \item \textbf{Deployment Optimization:} Apply model compression (pruning, quantization) and knowledge distillation to meet latency constraints, and develop calibrated uncertainty estimates for human-in-the-loop verification.
\end{enumerate}

% Context: Final summary emphasizing the practical value and direction of the work
\section{Closing Remarks}
The results underscore that robust deepfake detection benefits from integrating spatial, temporal, and cross-modal reasoning. By explicitly modeling audio--visual synchronicity and leveraging bidirectional temporal context, the proposed multimodal framework advances the frontier of practical and fair deepfake detection. Continued progress will require sustained attention to generalization, fairness, and adversarial resilience as generative methods evolve.